% The birth--plague model --- plan \S3.3, "\S6".
% Sources: 09/08 and 09/09.
% C4: 09/08:5,8 and 09/09:13,16 print "Corruption" and "Infection" with
%     identical left-hand sides.  Written here as one loss channel.
% C5: 09/09:11 introduces dV/dt but never couples V into a rate.  See the
%     NEEDS-AUTHOR marker in \cref{var:sec:vitality}.
% C6: 09/09:6 and 09/09:8 are the same paragraph twice, once rough and once
%     tidied.  The tidied one (:8) is kept.
% Renames: birth \beta -> \lambda (\beta is the speciation rate and nothing
% else; \lambda is the birth rate thesis-wide -- plan \S7.2 closing sentence);
% contact coefficient \alpha -> \chi and consumption p -> c (plan \S7.1).
% Cut: 09/09:25-76, "Plague in popular culture", parked.
% The four orphan sub-labels fig:figure1--fig:figure4 at 09/08:22,28,36,42 are
% replaced by one labelled float.

\section{The birth--plague model}
\label{var:sec:plague}

\subsection{The model}
\label{var:sec:plaguemodel}

The scarcity mechanism of the previous section drove catastrophe by a
potential $V$ carried alongside a single count. This model drives the loss
rate by a second count instead, so that what accelerates
the loss of one component is the size of the other. Let $X_t$ count
uninfected individuals and $Y_t$ the infected, and let the transitions in a
vanishing interval be
\begin{align}
\pr{\Delta X_t=1,\ \Delta Y_t=0\mid X_t,Y_t}
  &=\lambda X_t\,\Delta t+o(\Delta t),
  \label{var:eq:plague}\\
\pr{\Delta X_t=-1,\ \Delta Y_t=1\mid X_t,Y_t}
  &=\bigl(\delta+\chi Y_t\bigr)X_t\,\Delta t+o(\Delta t),\notag\\
\pr{\Delta X_t=0,\ \Delta Y_t=-1\mid X_t,Y_t}
  &=\mu Y_t\,\Delta t+o(\Delta t),\notag\\
\pr{\Delta X_t=\Delta Y_t=0\mid X_t,Y_t}
  &=1-\bigl[(\lambda+\delta+\chi Y_t)X_t+\mu Y_t\bigr]\Delta t+o(\Delta t).
  \notag
\end{align}

The middle channel is written once and carries two mechanisms superposed. An
uninfected individual is lost to \emph{corruption} spontaneously, at the
constant per-capita rate $\delta$, and to infection by contact with an
already infected individual, at rate $\chi$ per uninfected--infected pair.
Both produce the same jump, $X$ down by one and $Y$ up by one, so the only
thing that distinguishes them is how their rates scale: the first is linear
in $X_t$, the second bilinear in $X_t Y_t$. Writing them as separate lines
with the same left-hand side would state the same channel twice rather than
two channels once.

\begin{figure}[!htbp]
\centering
\begin{tikzpicture}[
  >={Stealth[length=1.9mm]},
  st/.style={draw=ink,thick,rounded corners=1.5pt,minimum width=15mm,
             minimum height=6.5mm,inner sep=1.5pt,font=\small,text=ink},
  lb/.style={font=\footnotesize,inner sep=1.8pt},
  every path/.style={thick}
]
\node[st,fill=type1!12] (c)  at (0,0)      {$(x,y)$};
\node[st]            (r)  at (4.4,0)    {$(x{+}1,y)$};
\node[st]            (dl) at (-4.4,-2.3) {$(x{-}1,y{+}1)$};
\node[st]            (d)  at (0,-2.3)   {$(x,y{-}1)$};
\draw[->] (c) -- node[above,lb]{$\lambda x$} (r);
\draw[->] (c) -- node[above left=-1pt,lb]{$(\delta+\chi y)\,x$} (dl);
\draw[->] (c) -- node[right,lb]{$\mu y$} (d);
\node[lb,anchor=west,text width=4.4cm] at (1.6,-2.3)
  {\itshape corruption and contact infection share one jump};
\end{tikzpicture}
\caption{Transitions out of the state $(x,y)$ in the birth--plague model,
drawn in the convention of \cref{m:fig:ratediagram}. Three channels leave
each interior state, and only the diagonal one carries a rate depending on
the other component --- which makes it the single place where feedback enters
the model: the more of the population is already infected, the faster the
rest is lost. Setting $\chi=0$ removes that dependence and leaves a linear
birth process for $X$ feeding an $\mathrm{M}/\mathrm{M}/\infty$-like queue,
with no coupling at all.}
\label{var:fig:plague-rates}
\end{figure}

Only one rate carries $Y_t$, but the coupling is genuinely two-way, because
every infection removes one $X$ and adds one $Y$: a large $Y_t$ both
accelerates the loss of $X$ and is replenished by that same loss. Running
against it are $\mu Y_t$, which clears infected individuals, and $\lambda
X_t$, which replaces uninfected ones. Whether the result settles or cycles is
a question about the balance of those three terms, and it is not answered
here.

\begin{figure}[!htbp]
\centering
\includegraphics[width=0.96\textwidth]{fig09_plague_paths/fig09}
\caption{Two realisations of \cref{var:eq:plague} at
$(\lambda,\delta,\chi,\mu)=(1,\,0.1,\,0.02,\,10)$: uninfected count $X_t$ in
blue, infected count $Y_t$ in vermillion. The clearance rate is large
relative to the contact rate, so infected individuals never accumulate and
the epidemic acts as a brake on $X_t$ rather than as a collapse of it. Set
beside \cref{var:fig:vbc-many}, the fluctuations never reset a phase.}
\label{var:fig:bpplots}
\end{figure}

\subsection{The vitality-coupled variant}
\label{var:sec:vitality}

The variant combines this model with that of \cref{var:sec:rrrr}, taking from
there the idea that catastrophe --- here the onset of plague rather than an
extinction --- is driven by scarcity per replicating agency, the agencies
being those thrown off by successive bifurcations. The shape is a loose one
and it fits a good many phenomena, most obviously the collapse of a
population that has outgrown what sustains it.

Carry a vitality $V(t)$ alongside the pair, replenished at a constant rate
and drawn down in proportion to the population:
\begin{equation}
\ddt{V}=\phi-c\,X_t .
\label{var:eq:vitalityV}
\end{equation}
The transitions are then those of \cref{var:eq:plague} with the spontaneous
part of the loss channel divided by the vitality:
\begin{equation}
\pr{\Delta X_t=-1,\ \Delta Y_t=1\mid X_t,Y_t}
=\lrb{\frac{\delta}{V(t)}+\chi Y_t}X_t\,\Delta t+o(\Delta t),
\label{var:eq:vitality}
\end{equation}
the other three channels being unchanged.

% NEEDS-AUTHOR: the source (09/09:11) specifies dV/dt = \xi - p X_t and then
% uses V in none of the transition rates below it, so the model as printed is
% \cref{var:eq:plague} plus a decorative extra equation.  The intent is stated
% in its prose: that plague instantiation is driven by relative per-capita
% scarcity, as in the rise--run--ruin--rejuvenation model.
% Written here as: corruption at rate \delta X_t / V(t), by consistency with
% the catastrophe channel \delta X_t / V(t) of \cref{var:eq:rrrr}, which is
% the model this variant says it is borrowing from.  The contact term \chi Y_t
% is left uncoupled, since contact infection is not a scarcity mechanism.
% Alternative reading: "per-capita scarcity" means scarcity per agency
% multiplied by the number of agencies exposed to it, giving
% \delta X_t^2 / V(t).  That is the reading under which the loss channel is
% quadratic in X and the collapse is sharper.
% A third possibility, that the whole loss channel including the contact term
% should be divided by V, is not adopted here for the reason given above.
% Author to confirm which was meant.  Nothing else in this chapter depends on
% the answer.

The reason for wanting the variant is that \cref{var:eq:plague} has no
mechanism for a collapse: its infected count rises when $X_t$ is large and
falls when it is small, which damps rather than amplifies. Dividing the
corruption rate by a vitality the population itself depletes supplies the
amplification that is missing, since a population which has grown large
drives $V$ down and so raises its own loss rate --- exactly the ingredient
\cref{var:sec:rrrr} used to close its cycle. Whether the combination gives a
sustained cycle or a single collapse followed by a low-level steady state is
the question the variant exists to ask, and it has not been simulated.

\paragraph{Open questions.}\label{var:rem:open-plague}
\begin{enumerate}[leftmargin=1.6em,itemsep=0.2em,topsep=0.2em]
\item \textbf{Extinction of $X$.} What is $\pr{X_t=0\text{ eventually}}$?
Not the standard birth--death calculation of \cref{m:sec:bd}: the loss rate
$(\delta+\chi Y_t)X_t$ depends on the second component, and the pair
$(X_t,Y_t)$ has no product-form generating function that closes. The obvious
route is a first-step analysis on the two-dimensional lattice, in the manner
of \cref{path:sec:rsfirststep}, and it has not been attempted.
\item \textbf{Cycles or a single collapse.} Does the vitality coupling of
\cref{var:eq:vitality} produce sustained cycles, and if so with what period?
This is the question \cref{var:sec:vitality} was written to pose rather than
to settle. Settling it wants simulation first and a first-passage analysis
second, and the answer will turn on whether corruption is $\delta X_t/V(t)$ as written,
or $\delta X_t^2/V(t)$.
\item \textbf{The uncoupled contact term.} If contact infection is left
outside the scarcity coupling, as it is in \cref{var:eq:vitality}, then at
large $V$ the model reduces to \cref{var:eq:plague} and at small $V$ it is
dominated by corruption. Whether that crossover is sharp enough to be
visible in a realisation is not known.
\end{enumerate}

\FloatBarrier
